% 1 Pagina (2 Max)
Actualmente, la detección del origen de las señales en electroencefalografía (EEG) sigue siendo un reto crítico para la investigación médica y clínica. Los métodos convencionales como \textit{sLORETA}, \textit{VARETA} y \textit{MNE}, aunque ampliamente utilizados, presentan limitaciones importantes: tienden a favorecer fuentes cercanas a la superficie cerebral, tienen baja resolución espacial, alta sensibilidad al ruido y requieren un alto costo computacional. Estas deficiencias impactan negativamente en la precisión diagnóstica, la investigación neurocientífica y el desarrollo de interfaces cerebro-computadora (BCI).

Ante esta necesidad, se propone un enfoque innovador que combine modelos inversos asistidos por aprendizaje profundo, inferencia bayesiana y herramientas de visualización interactiva. Esta integración permitirá identificar con mayor precisión las fuentes neuronales incluso en regiones profundas o con señales ruidosas, mejorando la confiabilidad de los análisis y facilitando su interpretación.

El desarrollo de esta herramienta tendrá múltiples beneficios: ofrecerá una solución abierta para el estudio del cerebro, mejorará la detección de patologías neurológicas como la epilepsia, y permitirá construir interfaces BCI más robustas para personas con movilidad reducida. Además, fomentará la adopción de estándares modernos de reproducibilidad, interpretabilidad y visualización neurofisiológica.
